\documentclass[11pt,a4paper]{article}
\usepackage{../shared/hanzocover}
\usepackage[utf8]{inputenc}
\usepackage[T1]{fontenc}
\usepackage{amsmath,amssymb}
\usepackage{graphicx}
\usepackage{booktabs}
\usepackage{hyperref}
\usepackage{xcolor}
\usepackage{listings}
\input{../shared/lstlang}
\usepackage{tikz}
\usetikzlibrary{shapes,arrows,positioning}

\lstset{
    basicstyle=\ttfamily\small,
    keywordstyle=\color{blue},
    commentstyle=\color{gray},
    stringstyle=\color{red},
    breaklines=true,
    frame=single,
    numbers=left,
    numberstyle=\tiny\color{gray}
}

\title{End-to-End Encryption and Security Architecture\\for Cloud-Native Applications}
\author{Antje Worring, Zach Kelling\thanks{zach@hanzo.ai} \\ \textit{CTO, Hanzo Industries (Techstars '17)}}
\date{2019}

\begin{document}

\hanzocoverpage

\begin{abstract}
We describe the encryption and security architecture developed at Hanzo Industries between 2018 and 2019, comprising four components: Age-based file encryption using X25519 key agreement (\texttt{hanzo/age}, May 2019), an encrypted message queue (\texttt{hanzo/mq}, December 2018, 4{,}900 commits), per-tenant storage encryption for the S3-compatible object store, and a key management architecture that serves as the precursor to the Hanzo KMS. Together, these components provide end-to-end encryption for data at rest, data in transit within the message queue, and data in storage. The security primitives described here subsequently evolved into the post-quantum encryption stack deployed on the Lux network.
\end{abstract}

\section{Introduction}

Cloud-native applications handle sensitive data across multiple services, storage backends, and communication channels. Encryption must be applied consistently at every layer: files at rest, messages in transit between services, objects in storage, and secrets in configuration. Relying on transport-layer encryption (TLS) alone is insufficient---data is decrypted at each service boundary, creating exposure windows.

The Hanzo security architecture addresses this through defense-in-depth encryption:

\begin{enumerate}
    \item \textbf{File encryption} (May 2019): Age-based encryption for configuration files, backups, and secrets at rest.
    \item \textbf{Message encryption} (December 2018): Per-message encryption within the message queue.
    \item \textbf{Storage encryption}: Per-tenant encryption of objects in the S3-compatible store.
    \item \textbf{Key management}: Centralized key lifecycle management.
\end{enumerate}

\section{Age Encryption}

\subsection{Design}

The \texttt{hanzo/age} module (May 2019) implements file encryption using the Age format~\cite{age2019}. Age uses X25519 key agreement with ChaCha20-Poly1305 symmetric encryption:

\begin{equation}
    K_{\text{shared}} = \text{X25519}(sk_{\text{sender}}, pk_{\text{recipient}})
\end{equation}

\begin{equation}
    K_{\text{file}} \xleftarrow{\$} \{0, 1\}^{256}
\end{equation}

\begin{equation}
    C_{\text{header}} = \text{ChaCha20-Poly1305}_{K_{\text{shared}}}(K_{\text{file}})
\end{equation}

\begin{equation}
    C_{\text{body}} = \text{ChaCha20-Poly1305}_{K_{\text{file}}}(P)
\end{equation}

where $sk_{\text{sender}}$ is the sender's private key, $pk_{\text{recipient}}$ is the recipient's public key, $K_{\text{file}}$ is a per-file symmetric key, $P$ is the plaintext, and $C_{\text{body}}$ is the ciphertext.

\subsection{Key Format}

Public keys are encoded as Bech32 strings with a human-readable prefix:

\begin{lstlisting}[caption=Age key pair format.]
# Public key
age1ql3z7hjy54pw3hyww5ayyfg7zqgvc7w3j2elw8zmrj2kg5sfn9aqmcac8p

# Private key (stored in KMS, never on disk)
AGE-SECRET-KEY-1QFWET6A35FH2XNNWCJLZ5VQCGY3P46S...
\end{lstlisting}

\subsection{Use Cases}

Age encryption is used for:
\begin{itemize}
    \item \textbf{Secret files}: Encrypting configuration files containing credentials before storing them in version control or backup storage.
    \item \textbf{Database backups}: Encrypting PostgreSQL and KV store backups before uploading to object storage.
    \item \textbf{Compliance exports}: Encrypting data exports for regulatory compliance.
\end{itemize}

\section{Encrypted Message Queue}

\subsection{Architecture}

The Hanzo message queue (\texttt{hanzo/mq}, December 2018, 4{,}900 commits) extends the PubSub layer with per-message encryption. Messages containing sensitive data are encrypted before publication and decrypted by authorized subscribers.

\subsection{Envelope Encryption}

Messages use envelope encryption: each message is encrypted with a unique data encryption key (DEK), and the DEK is encrypted with the recipient's key encryption key (KEK):

\begin{equation}
    DEK_m \xleftarrow{\$} \{0, 1\}^{256}
\end{equation}

\begin{equation}
    C_m = \text{AES-256-GCM}_{DEK_m}(M)
\end{equation}

\begin{equation}
    E_m = \text{AES-256-GCM}_{KEK_r}(DEK_m)
\end{equation}

where $M$ is the plaintext message, $C_m$ the ciphertext, $DEK_m$ the per-message data encryption key, and $KEK_r$ the recipient's key encryption key.

The published message contains $(E_m, \text{nonce}, C_m, \text{tag})$. Only subscribers with access to $KEK_r$ can decrypt $DEK_m$ and subsequently $M$.

\subsection{Key Rotation}

KEKs are rotated on a configurable schedule (default: 90 days). During rotation:
\begin{enumerate}
    \item A new KEK is generated and distributed to authorized subscribers.
    \item New messages are encrypted with the new KEK.
    \item The old KEK is retained for a grace period to allow in-flight message decryption.
    \item After the grace period, the old KEK is archived to the key management layer.
\end{enumerate}

\subsection{Message Queue Performance}

\begin{table}[h]
\centering
\caption{Encryption overhead on message throughput.}
\begin{tabular}{lrr}
\toprule
\textbf{Mode} & \textbf{Messages/s} & \textbf{Overhead} \\
\midrule
Plaintext PubSub   & 3{,}500{,}000 & --- \\
Encrypted MQ (256B) & 2{,}800{,}000 & 20\% \\
Encrypted MQ (4KB)  & 1{,}900{,}000 & 46\% \\
Encrypted MQ (64KB) &   420{,}000 & 88\% \\
\bottomrule
\end{tabular}
\end{table}

The overhead is dominated by the symmetric encryption of the message body. The envelope key operation adds negligible latency (sub-microsecond).

\section{Per-Tenant Storage Encryption}

\subsection{Architecture}

The S3-compatible object storage layer encrypts all objects at rest using per-tenant keys. The encryption is transparent to clients: objects are encrypted on write and decrypted on read by the storage service.

\subsection{Key Hierarchy}

A three-tier key hierarchy provides isolation and rotation:

\begin{enumerate}
    \item \textbf{Master Key (MK)}: A single master key stored in a hardware security module (HSM) or software KMS. Never leaves the key management boundary.
    \item \textbf{Tenant Key (TK)}: One per tenant, encrypted by the MK. Stored alongside the tenant metadata.
    \item \textbf{Object Key (OK)}: One per object, encrypted by the TK. Stored in the object metadata.
\end{enumerate}

\begin{equation}
    OK_o \xleftarrow{\$} \{0, 1\}^{256}, \quad C_o = \text{AES-256-GCM}_{OK_o}(O)
\end{equation}

\begin{equation}
    E_{OK} = \text{AES-256-GCM}_{TK_t}(OK_o), \quad E_{TK} = \text{AES-256-GCM}_{MK}(TK_t)
\end{equation}

\subsection{Tenant Isolation Guarantee}

Compromising a single tenant key $TK_t$ exposes only that tenant's objects. The master key and all other tenant keys remain unaffected. Tenant key rotation re-encrypts only the tenant key itself (not the underlying objects), making rotation a constant-time operation:

\begin{equation}
    E_{TK}' = \text{AES-256-GCM}_{MK}(TK_t'), \quad \forall\, o : E_{OK_o}' = \text{AES-256-GCM}_{TK_t'}(OK_o)
\end{equation}

In practice, lazy re-encryption of object keys is performed during object access, amortizing the cost over normal read/write operations.

\section{Key Management Architecture}

\subsection{Design Principles}

The key management layer follows three principles:

\begin{enumerate}
    \item \textbf{Separation of duties}: The system that encrypts data does not have direct access to the master key. Encryption operations go through the KMS API.
    \item \textbf{Audit trail}: Every key operation (creation, rotation, access, deletion) is logged immutably.
    \item \textbf{Least privilege}: Services are granted access only to the specific keys required for their function.
\end{enumerate}

\subsection{KMS API}

The key management API provides four operations:

\begin{lstlisting}[language=Go,caption=Key management interface.]
type KMS interface {
    // Create a new encryption key
    CreateKey(ctx context.Context,
        name string) (*Key, error)

    // Encrypt data with a named key
    Encrypt(ctx context.Context,
        keyName string,
        plaintext []byte) ([]byte, error)

    // Decrypt data with a named key
    Decrypt(ctx context.Context,
        keyName string,
        ciphertext []byte) ([]byte, error)

    // Rotate a key (creates new version)
    RotateKey(ctx context.Context,
        keyName string) (*Key, error)
}
\end{lstlisting}

\subsection{Access Control}

Access to keys is controlled by IAM policies. Each key has an access control list specifying which service identities may encrypt, decrypt, or administer the key:

\begin{table}[h]
\centering
\caption{Key access policy example.}
\begin{tabular}{lll}
\toprule
\textbf{Service} & \textbf{Key} & \textbf{Permissions} \\
\midrule
storage-service & tenant-keys/* & encrypt, decrypt \\
backup-service  & backup-key    & encrypt only \\
restore-service & backup-key    & decrypt only \\
kms-admin       & *             & admin \\
\bottomrule
\end{tabular}
\end{table}

\section{Threat Model}

The encryption architecture defends against the following threats:

\begin{itemize}
    \item \textbf{Storage breach}: An attacker gaining read access to the storage backend obtains only ciphertext. Without the tenant key, objects cannot be decrypted.
    \item \textbf{Message interception}: An attacker capturing messages from the queue obtains only encrypted payloads. Without the recipient KEK, messages cannot be decrypted.
    \item \textbf{Insider threat}: A compromised service can only access keys granted to its identity. Separation of encrypt/decrypt permissions limits blast radius.
    \item \textbf{Key compromise}: Per-tenant and per-object key hierarchy limits the scope of any single key compromise.
\end{itemize}

\section{Evolution Toward Post-Quantum Security}

The X25519 key agreement and AES-256-GCM symmetric encryption provide strong security against classical adversaries. However, the emergence of large-scale quantum computing threatens the security of elliptic curve key agreement. The key management architecture was designed with algorithm agility: the KMS abstraction allows replacement of the underlying cryptographic primitives without changing the API surface.

This design decision enabled the subsequent adoption of post-quantum key encapsulation mechanisms (KEMs) on the Lux network, where Kyber/ML-KEM replaced X25519 for key agreement while maintaining the same envelope encryption pattern and key hierarchy.

\section{Conclusion}

The Hanzo encryption and security architecture provides defense-in-depth protection for data at rest, in transit within the message queue, and in object storage. Age encryption handles file-level encryption with X25519 key agreement. The encrypted message queue extends PubSub with per-message envelope encryption. Per-tenant storage encryption with a three-tier key hierarchy ensures tenant isolation. The key management API centralizes key lifecycle operations with IAM-controlled access policies. These primitives, developed between 2018 and 2019, formed the security foundation that subsequently evolved into the post-quantum encryption stack deployed on the Lux network.

\bibliographystyle{plain}
\begin{thebibliography}{9}

\bibitem{age2019}
F.~Valsorda. age: A Simple, Modern, and Secure Encryption Tool. \url{https://age-encryption.org}, 2019.

\bibitem{rfc7748}
A.~Langley, M.~Hamburg, and S.~Turner. Elliptic Curves for Security. RFC 7748, IETF, 2016.

\bibitem{rfc5116}
D.~McGrew. An Interface and Algorithms for Authenticated Encryption. RFC 5116, IETF, 2008.

\bibitem{nist2016}
NIST. Recommendation for Block Cipher Modes of Operation: Galois/Counter Mode (GCM). SP 800-38D, 2007.

\bibitem{bernstein2012}
D.~J.~Bernstein. Curve25519: New Diffie-Hellman speed records. \textit{Lecture Notes in Computer Science}, 3958:207--228, 2006.

\end{thebibliography}

\end{document}
